% --- Archon physics formalization source begin ---
% archon:physics
% archon:covers PhyXMiniProblems/problem_phyx_mini_0005.lean
% archon:source-report reports/phyx_mini/problem_phyx_mini_0005.source.json

\chapter{Physics problem phyx\_mini\_0005}
\label{ch:PhyXMiniProblems_problem_phyx_mini_0005}

\paragraph{Problem source.}
\#\# Physical scenario

A person looking into an empty container is able to see the far edge of the container's bottom as shown in figure. The height of the container is \textbackslash{}( h \textbackslash{}), and its width is \textbackslash{}( d \textbackslash{}). When the container is completely filled with a fluid of index of refraction \textbackslash{}( n \textbackslash{}) and viewed from the same angle, the person can see the center of a coin at the middle of the container’s bottom

\#\# Question

For what range of values of \textbackslash{}( n \textbackslash{}) will the center of the coin not be visible for any values of \textbackslash{}( h \textbackslash{}) and \textbackslash{}( d \textbackslash{})?

\#\# Answer choices

- A: A: \textbackslash{}( n>2.5 \textbackslash{})

- B: B: \textbackslash{}( n>1.5 \textbackslash{})

- C: C: \textbackslash{}( n>2 \textbackslash{})

- D: D: \textbackslash{}( n>5 \textbackslash{})

\#\# Figure caption (auxiliary; use the image as primary evidence)

The image shows a vertical cross-section of a transparent container, such as a glass, with a coin at the bottom. 

Key elements include:

1. **Container**: It has slightly curved sides and is open at the top.
2. **Coin**: Located at the bottom center of the container.
3. **Eye**: Positioned above and to the side of the container, observing the coin.
4. **Arrow**: A blue arrow represents the line of sight from the eye to the coin.
5. **Text and Labels**:
   - The height from the top of the container to the coin is labeled as \textbackslash{}( h \textbackslash{}).
   - The horizontal distance from the side of the container to the point where the line of sight enters the container is labeled as \textbackslash{}( d \textbackslash{}).

The diagram likely illustrates a concept related to optics or refraction.

\#\# Dataset metadata

Geometrical Optics; Physical Model Grounding Reasoning, Spatial Relation Reasoning

\paragraph{Recorded answer/context.}
C: C: \textbackslash{}( n>2 \textbackslash{})

\paragraph{Figure/image path.}
/root/proposal\_for\_physic/science-mango/phyx\_mini\_run/phyx\_data/test\_image/5.png

\paragraph{Formalization target.}
create a compiling Lean file with sorry bodies at `PhyXMiniProblems/problem\_phyx\_mini\_0005.lean`. The Lean declarations must preserve the physical quantities, dimensions or dimensional roles, figure labels, governing-law hypotheses, and final relation expressed by this problem.
Use Mathlib/Physlib names found through LeanExplore where available. If a physics API is missing, introduce faithful local abstractions rather than scalar placeholder aliases.

\begin{theorem}[Physics formalization target]
\label{thm:physics:phyx_mini_0005:target}
\lean{PhyXMiniProblems.ProblemPhyXMini0005.refractiveIndex_gt_two_makes_coinCenter_invisible}
\uses{def:physics:phyx-mini-0005:phyxminiproblems-problemphyxmini0005-containergeometry, def:physics:phyx-mini-0005:phyxminiproblems-problemphyxmini0005-emptycontainerfaredgesightline, def:physics:phyx-mini-0005:phyxminiproblems-problemphyxmini0005-coincentervisibleatsameviewingangle}
If the fluid refractive index is greater than \(2\), then for every positive
container height and width the coin centre cannot be visible at the same
external angle from which the far bottom edge was visible in the empty
container.
\end{theorem}
\begin{proof}
Write \(x=d/h\), which is positive.  The empty-container sightline gives
\(\tan\alpha=x\), whereas visibility of the bottom midpoint would give an
acute fluid angle \(\beta\) with \(\tan\beta=x/2\).  For an acute angle
\(\theta\),
\[
  \sin\theta=\frac{\tan\theta}{\sqrt{1+\tan^2\theta}}.
\]
Consequently
\[
 \frac{\sin\alpha}{\sin\beta}
 =\frac{\sqrt{4+x^2}}{\sqrt{1+x^2}}<2,
\]
where the strict inequality follows after squaring from \(x>0\).  Both sines
are positive.  Snell's law would therefore force the fluid refractive index
to equal this ratio and hence be less than \(2\), contradicting the assumed
strict lower bound.
\end{proof}

% --- Archon declaration topology begin ---
\section{Declaration topology}
\begin{definition}[si Length Value]
  \label{def:physics:phyx-mini-0005:phyxminiproblems-problemphyxmini0005-silengthvalue}
  \lean{PhyXMiniProblems.ProblemPhyXMini0005.siLengthValue}
  The real-valued readout of a physical length when it is expressed in SI units
\end{definition}
\begin{proof}
  This is the declared model definition of si Length Value.
\end{proof}
\begin{definition}[Container Geometry]
  \label{def:physics:phyx-mini-0005:phyxminiproblems-problemphyxmini0005-containergeometry}
  \lean{PhyXMiniProblems.ProblemPhyXMini0005.ContainerGeometry}
  \uses{def:physics:phyx-mini-0005:phyxminiproblems-problemphyxmini0005-silengthvalue}
  The dimensionful labels in the container cross-section. The coin is at the midpoint of the bottom segment of length width; height is the vertical distance from the bottom to the horizontal fluid surface and viewing rim
\end{definition}
\begin{proof}
  This is the declared model definition of Container Geometry.
\end{proof}
\begin{definition}[Snell Law At Fluid Air Interface]
  \label{def:physics:phyx-mini-0005:phyxminiproblems-problemphyxmini0005-snelllawatfluidairinterface}
  \lean{PhyXMiniProblems.ProblemPhyXMini0005.SnellLawAtFluidAirInterface}
  Snell's law for a ray leaving a fluid through its horizontal surface into air of refractive index one. The refractive index is dimensionless, and both angle readouts are real numbers measured in radians from the surface normal
\end{definition}
\begin{proof}
  This is the declared model definition of Snell Law At Fluid Air Interface.
\end{proof}
\begin{definition}[Empty Container Far Edge Sightline]
  \label{def:physics:phyx-mini-0005:phyxminiproblems-problemphyxmini0005-emptycontainerfaredgesightline}
  \lean{PhyXMiniProblems.ProblemPhyXMini0005.EmptyContainerFarEdgeSightline}
  \uses{def:physics:phyx-mini-0005:phyxminiproblems-problemphyxmini0005-silengthvalue, def:physics:phyx-mini-0005:phyxminiproblems-problemphyxmini0005-containergeometry}
  The empty-container figure readout at the fixed viewing angle: the sightline through the near top rim reaches the far bottom edge, so its horizontal run is the full width d and its vertical run is the height h
\end{definition}
\begin{proof}
  This is the declared model definition of Empty Container Far Edge Sightline.
\end{proof}
\begin{definition}[Coin Center Visible At Same Viewing Angle]
  \label{def:physics:phyx-mini-0005:phyxminiproblems-problemphyxmini0005-coincentervisibleatsameviewingangle}
  \lean{PhyXMiniProblems.ProblemPhyXMini0005.CoinCenterVisibleAtSameViewingAngle}
  \uses{def:physics:phyx-mini-0005:phyxminiproblems-problemphyxmini0005-silengthvalue, def:physics:phyx-mini-0005:phyxminiproblems-problemphyxmini0005-containergeometry, def:physics:phyx-mini-0005:phyxminiproblems-problemphyxmini0005-snelllawatfluidairinterface}
  The center of the coin is visible at the same external viewing angle when an acute ray from the bottom midpoint to the near top rim has the figure-imposed half-width slope and obeys Snell's law at the fluid--air interface
\end{definition}
\begin{proof}
  This is the declared model definition of Coin Center Visible At Same Viewing Angle.
\end{proof}
% --- Archon declaration topology end ---
% --- Archon physics formalization source end ---
